% 第 19 讲 布劳尔群（源：src/lecture19.txt；已删去抽出时残留的页眉「19. BRAUER GROUPS 83/85」）
\section*{第 19 讲\quad 布劳尔群}

本讲我们概述布劳尔群。叙述将略于证明，但会给出精确的构造与命题陈述；完整证明见 [Mil13, Ser79, Boy07]。我们的动机问题是：「那些关于哈密顿代数的东西究竟是怎么回事？」（见习题集 1 第 5 题、习题集 2 第 3 题与习题集 3 第 2 题。）我们将看到，有两个对象都叫作「布劳尔群」：一个由上同调定义，另一个是更一般的代数定义；我们将证明二者一致。

回忆：若 $L/K$ 是非阿基米德局部域的 $G$-Galois 扩张，则
\[
(\Z[-2])^{tG}\simeq (L^\times)^{tG}.
\]
取 $H^0$（零次上同调）时，我们得到 lcft，即
\[
K^\times/NL^\times\simeq H^1(G,\Z)=G^{\mathrm{ab}}.
\]
这很自然，因为 $L^\times$ 位于 $0$ 次。但 $\Z$ 位于 $2$ 次，那么取 $H^2$ 会怎样？此时我们得到同构
\[
\widehat H^0(G,\Z)\simeq\widehat H^2(G,L^\times),
\]
其左端极简单地同构于
\[
\Z/[L:K]\Z=\Z/\#G\Z=\tfrac{1}{[L:K]}\Z/\Z,
\]
因为 $G$ 平凡地作用在 $\Z$ 上，故不变元正是 $\Z$，而范数对应乘以 $\#G$。然而右端要神秘得多，这促使下述定义。

\textbf{定义 19.1.} $L/K$ 的上同调布劳尔群为
\[
\Br^{\mathrm{coh}}(K/L):=H^2(G,L^\times).
\]

\textbf{注 19.2.} 若变动 $L$ 会怎样？设 $L_2/L_1/K$ 为 Galois 扩张，$\Gal(L_i/K)=G_i$（$i=1,2$）。则有短正合列
\[
0\to \Gal(L_2/L_1)\to G_2\to G_1\to 0,
\]
以及映射
\[
\Br^{\mathrm{coh}}(L_1/K)=H^2(G_1,L_1^\times)\to H^2(G_2,L_1^\times)\to H^2(G_2,L_2^\times)=\Br^{\mathrm{coh}}(L_2/K),
\]
因为关于 $G_2$ 的不变性蕴含关于 $G_1$ 的不变性，并且经由嵌入 $L_1^\times\hookrightarrow L_2^\times$ 过渡。这促使下述定义。
% 原文此处写 $\Br^{\mathrm{coh}}(L_1/K)$，而定义 19.1 写的是 $\Br^{\mathrm{coh}}(K/L)$；此处照原文排，未统一。

\textbf{定义 19.3.} $K$ 的上同调布劳尔群为
\[
\Br^{\mathrm{coh}}(K):=\varinjlim_{L/K}\Br^{\mathrm{coh}}(K/L)=H^2(\Gal(\overline K/K),\overline K^\times),
\]
其中正向极限取遍有限 Galois 扩张 $L/K$。

注意上式最右端的表达用了我们上一讲的记号。

\textbf{断言 19.4.} 在 lcft 之下，下述图表交换：
% [按原 PDF 第 86 页（印刷页 82）复原：上排 $H^2(G_2,L_1^\times)\to H^2(G_1,L_1^\times)$、下排 $\frac{1}{[L_1:K]}\Z/\Z \leftarrow \frac{1}{[L_2:K]}\Z/\Z$（左向包含），两列竖直箭头均为「$\|$」，即 lcft 给出的同构]
\[
\begin{CD}
H^2(G_2,L_1^\times) @>>> H^2(G_1,L_1^\times)\\
@V{\sim}VV @V{\sim}VV\\
\tfrac{1}{[L_1:K]}\Z/\Z @<<< \tfrac{1}{[L_2:K]}\Z/\Z
\end{CD}
\]
其中竖直的同构即 lcft，下方水平箭头为自然的包含 $\tfrac{1}{[L_2:K]}\Z/\Z\hookrightarrow\tfrac{1}{[L_1:K]}\Z/\Z$。

\textbf{推论 19.5.} 对非阿基米德局部域 $K$，有
\[
\Br^{\mathrm{coh}}(K)\simeq\Q/\Z.
\]

\textbf{注 19.6.} 还可以证明最上方的映射是单射。对非阿基米德局部域的扩张 $L/K$，有正合列
\[
0\to \Br^{\mathrm{coh}}(K/L)\to
\underbrace{\Br^{\mathrm{coh}}(K)}_{H^2(\Gal(\overline K/K),\overline K^\times)}\to
\underbrace{\Br^{\mathrm{coh}}(L)}_{H^2(\Gal(\overline K/L),\overline K^\times)},
\]
我们下次课将说明其理由。

现在转向代数视角，它给出布劳尔群的经典定义。

\textbf{命题 19.7.} 设 $K$ 为域，$A$ 是有限维 $K$-代数，其中心为 $K$。则下述条件等价：
\begin{enumerate}[label=(\arabic*)]
\item $A$ 是单的，即它没有非平凡的双边理想。
\item 对某个可分扩张 $L/K$，$A\otimes_K L\simeq M_n(L)$（即 $n\times n$ 矩阵代数）。
\item $A\simeq M_n(D)$，其中 $D$ 是 $K$ 上的中心（即以 $K$ 为中心）可除（即乘法逆元存在，但乘法未必交换）代数。
\end{enumerate}
若这些条件成立，则称 $A$ 为 $K$ 上的\textbf{中心单代数}（central simple algebra，csa；亦称 Azumaya 代数）。

\textbf{推论 19.8.} $K$ 上任何中心单代数的维数都是平方数。

\textit{证明.} $A$ 的维数在 (2) 中的张量运算下保持不变，而 $M_n(L)$ 的维数是平方数。\hfill$\square$

\textbf{例 19.9.} (1) $M_n(K)$。
(2) $K$ 上的中心可除代数。
(3) $\R$ 上的哈密顿（四元数）代数，此时 $K=\R$。
(4) 对一切满足 $\operatorname{char}(K)\neq 2$ 的域 $K$ 与元素 $a,b\in K^\times$，$H_{a,b}$ 是中心单代数。

然而，一般的域扩张 $L/K$（把 $L$ 看作 $K$-代数）并不具有中心性，其维数也未必是平方数。

\textbf{定义 19.10.} $K$ 的中心单代数布劳尔群为
\[
\Br^{\mathrm{csa}}(K):=\{\,K\ \text{上中心单代数的等价类}\,\},
\]
其中的等价关系为：$A\sim B$ 当且仅当 $A$ 与 $B$ 是上述 (3) 中同构的可除代数上的矩阵代数（维数不必相同）。

\textbf{注 19.11.} $A$ 的同构类只依赖于 $D$ 的同构类。
% 原文如此（「同构类」）；按上下文似应为「等价类」，照译并登记于报告。

\textbf{命题 19.12.} 若 $A$ 与 $B$ 是 $K$ 上的中心单代数，则 $A\otimes_K B$ 也是中心单代数（注意张量积使维数相乘，故 $A\otimes_K B$ 的维数也是平方数）。模等价类，这仅依赖于 $[A],[B]\in\Br^{\mathrm{csa}}(K)$，因此 $\Br^{\mathrm{csa}}(K)$ 构成一个阿贝尔群。

证明从略；证明 $A\otimes_K B$ 是单的相当麻烦。但我们指出，$K$ 上的矩阵代数代表 $\Br^{\mathrm{csa}}(K)$ 的单位元（或「$0$ 等价类」）。$A$ 的逆是 $A$ 但取反乘法（即 $x\cdot y:=yx$），我们记作 $A^{\mathrm{op}}$。事实上，有典范的代数同态
\[
A\otimes_K A^{\mathrm{op}}\to \operatorname{End}_K(A),
\]
其中 $A$ 与 $A^{\mathrm{op}}$ 从相反的两侧作用（注意 $\operatorname{End}_K(A)$ 属于 $\Br^{\mathrm{csa}}(K)$ 的 $0$ 等价类）。这个（非零）映射的核必是双边理想，而 $A\otimes_K A^{\mathrm{op}}$ 是单的，故该映射必是单射。又因两端在 $K$ 上的维数都是 $\dim_K(A)^2$，它因而也是满射，从而是同构。

\textbf{定义 19.13.} $L/K$ 的中心单代数布劳尔群为
\[
\Br^{\mathrm{csa}}(K/L):=\{\,\text{满足 } A\otimes_K L\simeq M_n(L)\ \text{的中心单代数 } A\ \text{的等价类}\,\},
\]
并称这样的 $A$ 在 $L$ 上\textbf{分裂}。

这一概念等价于底层的可除代数 $D$ 在 $L$ 上分裂，后者同样意味着 $D\otimes_K L\simeq M_n(L)$ 不再是可除代数。于是显然
\[
\Br^{\mathrm{csa}}(K)=\bigcup_{\substack{L/K\\ \text{可分}}}\Br^{\mathrm{csa}}(K/L).
\]

\textbf{例 19.14.} (1) 若 $K$ 代数闭，则 $\Br^{\mathrm{csa}}(K)=\{0\}$ 显然平凡。
(2) $\Br^{\mathrm{csa}}(\R)=\{\R,\mathbb H\}=\Z/2\Z$，因为 $\R$ 上唯一的另一个可除代数是 $\C$，而它不是中心的。注意哈密顿四元数代数 $\mathbb H$ 在 $\C$ 上分裂，维数为 $4=2^2$。
(3) $H_{a,b}$ 在 $K(\sqrt a)$ 与 $K(\sqrt b)$ 上分裂，它们正是 $H_{a,b}$ 中通常的交换子代数。事实上，对任意 $c\in K$ 都可取 $K(\sqrt c)$，还可以用单位作共轭。因此在下述断言中，$L$ 非常不唯一。注意与这一情形不同，这样的 $L$ 不必是 $K$ 的 Galois 扩张；例子很难找，但 $1970$ 年代已被发现。

\textbf{断言 19.15.} $n^2$ 维中心可除代数 $D/K$ 在 $n$ 次扩张 $L/K$ 上分裂，当且仅当 $K\subseteq L\subseteq D$ 作为子代数。等价地，$L$ 是 $D$ 中的极大交换子代数。

首先有一个（简单的）事实：可除代数的任何交换子代数都是域。这类似于「域上任何有限维整环本身是域」。这是因为乘以任何元素都是单射，且由矩阵给出，故是满射；于是某个元素映到单位元，给出所需的逆元。其次有一个（不显然的）事实：$\dim_K(D)$ 是 $D$ 在 $K$ 上任何极大交换子代数维数的平方。我们来证明断言的一个方向：

\textit{证明.} 设 $K\subseteq L\subseteq D$，$[L:K]=n$，$\dim_K(D)=n^2$。我们断言 $D$ 在 $L$ 上分裂。有映射
\[
D\otimes_K L\to \operatorname{End}_L(D)=M_n(L),
\]
其中 $L$ 经右乘作用在 $D$ 上，$D$ 经左乘作用在自身上；两个作用交换。该映射是单射（若对某个 $d\in D$、$l\in L$ 及所有 $x\in D$ 有 $dxl=x$，则 $dl=1$，从而对所有 $x\in D$ 有 $dx=xl^{-1}=xd$，故 $d\in K$，于是 $d\otimes l=1\otimes 1$），因此是满射，从而是同构，这正是所要。\hfill$\square$

\textbf{定理 19.16.} 对每个 Galois 扩张 $L/K$，有 $\Br^{\mathrm{coh}}(K)\simeq\Br^{\mathrm{csa}}(K)$ 以及 $\Br^{\mathrm{coh}}(K/L)\simeq\Br^{\mathrm{csa}}(K/L)$。

注意上同调布劳尔群只对 Galois 扩张 $L/K$ 定义，而中心单代数布劳尔群对所有扩张 $L/K$ 定义。这使布劳尔群的上同调定义有了意义。我们给出下述（不完整的）证明概要：

\textit{证明.} 对 $G$-Galois 扩张 $L/K$，定义映射
\begin{equation*}
H^2(G,L^\times)\longrightarrow \Br^{\mathrm{csa}}(K/L). \tag{19.1}
\end{equation*}
$H^2$ 中每个元素都由一个 $2$-上闭链代表，即满足
\[
g_1\phi(g_2,g_3)\phi(g_1,g_2g_3)=\phi(g_1g_2,g_3)\phi(g_1,g_2)
\]
的映射 $\phi:G\times G\to L^\times$（对每个 $g_1,g_2,g_3\in G$），其中我们引入了记号 ${}^g x:=g\cdot x$。对每个这样的 $\phi$，如下构造一个中心单代数：作
\[
L[G]=\Bigl\{\sum_{g\in G}x_g[g]:x_g\in L\Bigr\}
\]
并附以关系
\[
[g]x={}^g x[g],\qquad [g][h]=\phi(g,h)[gh]
\]
（对每个 $x\in L$ 与 $g,h\in G$）。下面验证 $2$-上闭链恒等式等价于结合律：
\begin{align*}
[g_1]([g_2][g_3]) &= [g_1]\phi(g_2,g_3)[g_2g_3]\\
&= {}^{g_1}\phi(g_2,g_3)[g_1][g_2g_3]\\
&= {}^{g_1}\phi(g_2,g_3)\phi(g_1,g_2g_3)[g_1g_2g_3]\\
&= \phi(g_1g_2,g_3)\phi(g_1,g_2)[g_1g_2g_3]\\
&= \phi(g_1,g_2)[g_1g_2][g_3]\\
&= ([g_1][g_2])[g_3],
\end{align*}
对一切 $g_1,g_2,g_3\in G$ 成立。我们断言（但不证明）这个中心单代数的等价类只依赖于 $\phi$ 模上边界，并且它在 $L$ 上分裂。此外，映射 (19.1) 是群同构。\hfill$\square$

这并非一个特别深刻的定理，尽管它远非显然；完整证明用到了大量并不特别好记的结构理论。

我们现在问：$K$ 上次数为 $n^2$ 的中心可除代数有多少个同构类？当 $n=1$ 时只有一个；当 $n=2$ 时同样只有一个，如习题集 3 第 2(d) 题所示。
% [原注：原文作 "When n = 1^2, there is only 1; when n = 2^2, there is again only 1"，即按次数 $n^2$（$n=1,2$，次数 1 与 4）表述；译文按 $n=1,2$ 排，语义一致（已按原 PDF 第 84 页核实，非抽出误字）]

\textbf{断言 19.17.} $K$ 上次数为 $n^2$ 的中心可除代数有 $\varphi(n)=\#(\Z/n\Z)^\times$ 个同构类。

\textbf{断言 19.18.} 设 $L/K$ 是非阿基米德局部域的 $n$ 次扩张。则下述图表交换：
\[
\begin{CD}
\Br(K) @>>> \Br(L)\\
@V{\mathrm{lcft}}VV @VV{\mathrm{lcft}}V\\
\Q/\Z @>{\times n}>> \Q/\Z
\end{CD}
\]

注意：鉴于定理 19.16，我们干脆把布劳尔群记作 $\Br$。事实上，我们已经看到两个水平映射都应有非平凡的核；当 $L=\overline K$ 时，$\Br(L)$ 平凡。于是次数为 $n^2$ 的中心可除代数在 $\Br(K)$ 中的阶为 $n$，即它是 $n$-挠的，因为它在自己的 $n$ 次极大交换子代数上分裂。等价地，$H^2(G,-)$ 是 $\#G$-挠的，如习题集 7 第 2(c) 题所证。

\textbf{推论 19.19.} 任何次数为 $n^2$ 的可除代数在 $K$ 的任何 $n$ 次扩张上分裂。

另一方面，若 $A\simeq M_n(D)$ 是 $K$ 上次数为 $n^2$ 的中心单代数，代表 $\Br(K)$ 中某个 $n$-挠类，则 $A$ 在 $D$ 的极大交换子代数上分裂。这意味着：若 $A$ 是 $n$-挠的，但对任何 $m\mid n$ 都不是 $m$-挠的，则 $A$ 本身是可除代数。这给出 $\varphi(n)$：可除代数来自 $\Z/n\Z$ 中那些不来自更小的 $\Z/m\Z$ 的类。另一个推论如下：

\textbf{推论 19.20.} $K$ 上任何次数为 $n^2$ 的中心可除代数都包含 $K$ 的每个 $n$ 次域扩张。
% 原文抽出作 "every degree-n field extension of L"，按上下文应为「$K$ 的」；照上下文译，并登记于报告。

\textit{证明.} 由断言 19.18，$\Br(K)$ 中任何 $n$-挠类在 $K$ 的任何 $n$ 次扩张 $L$ 上的布劳尔群 $\Br(L)$ 中像为零。故它在 $L$ 上分裂，从而由断言 19.15 包含 $L$。\hfill$\square$

对 $n=2$，这就是「每个 $x\in K$ 在 $H_{a,b}$ 中有平方根」这一结果，它已在习题集 3 第 2(c) 题中证明。
